% Sources: Eguchi--Ooguri--Tachikawa 2011, arXiv:1004.0956;
%   Eguchi--Hikami 2010, arXiv:1008.4924, Tables 1 and 3;
%   Gaberdiel--Hohenegger--Volpato 2010, arXiv:1008.3778, Tables 2 and 3;
%   Gaberdiel--Persson--Ronellenfitsch--Volpato 2013, arXiv:1211.7074;
%   Conway--Sloane ATLAS of Finite Groups, Mathieu $M_{24}$.

\providecommand{\kBKM}{\kappa_{\mathrm{BKM}}}
\providecommand{\CM}{C_{M_{24}}}

\section*{The 26 conjugacy classes of $M_{24}$ and twined K3
elliptic-genus numerics}

\paragraph{Conventions.} ATLAS labels are retained; rows are ordered by
element order. The frame shape is the cycle decomposition of \(g\) in the
natural \(24\)-point permutation representation. Let
\[
  Z_g(\tau,z)=\sum_{n,\ell}c_g(n,\ell)\,q^n y^\ell
\]
be the full K3 twining genus. Its polar coefficients are
\(c_g(0,1)=c_g(0,-1)=2\), and its Witten index is
\[
  \chi_{24}(g)=\operatorname{Tr}_{V_{24}}(g)
  =c_g(0,0)+4 .
\]
The columns \(A_g(1)\) and \(A_g(2)\) are the first two massive
\(\mathcal N=4\) traces:
\[
  A_g(1)=\operatorname{Tr}_{45+\overline{45}}(g),\qquad
  A_g(2)=\operatorname{Tr}_{231+\overline{231}}(g).
\]
The CHL column is populated exactly for the five primitive denominator
rows \(N\in\{1,2,3,4,6\}\).

\paragraph{The canonical 26-class table.}

{\footnotesize
\begin{center}
\begin{tabular}{l|l|c|r|r|r|r|r|c|c}
\toprule
 [g] & frame shape & $|g|$ & $|\CM(g)|$ &
   $\chi_{24}(g)$ & $c_g(0,0)$ &
   $A_g(1)$ & $A_g(2)$ & anom.\ & CHL $N$ \\
\midrule
 1A  & $1^{24}$                       & 1  & $244{,}823{,}040$ & $24$ & $20$ & $\phantom{-}90$ & $\phantom{-}462$ & no  & 1 \\
 2A  & $1^{8}2^{8}$                   & 2  & $21{,}504$        & $\phantom{0}8$ & $\phantom{-}4$  & $-6$ & $\phantom{-}14$  & no  & 2 \\
 2B  & $2^{12}$                       & 2  & $7{,}680$         & $\phantom{0}0$ & $-4$ & $\phantom{-}10$ & $-18$ & no  & -- \\
 3A  & $1^{6}3^{6}$                   & 3  & $1{,}080$         & $\phantom{0}6$ & $\phantom{-}2$  & $\phantom{-}0$  & $-6$  & no  & 3 \\
 3B  & $3^{8}$                        & 3  & $504$             & $\phantom{0}0$ & $-4$ & $\phantom{-}6$  & $\phantom{-}0$  & no  & -- \\
 4A  & $2^{4}4^{4}$                   & 4  & $384$             & $\phantom{0}0$ & $-4$ & $-6$ & $-2$ & no  & -- \\
 4B  & $1^{4}2^{2}4^{4}$              & 4  & $128$             & $\phantom{0}4$ & $\phantom{-}0$  & $\phantom{-}2$  & $-2$ & no  & 4 \\
 4C  & $4^{6}$                        & 4  & $96$              & $\phantom{0}0$ & $-4$ & $\phantom{-}2$  & $\phantom{-}6$  & no  & -- \\
 5A  & $1^{4}5^{4}$                   & 5  & $60$              & $\phantom{0}4$ & $\phantom{-}0$  & $\phantom{-}0$  & $\phantom{-}2$  & no  & -- \\
 6A  & $1^{2}2^{2}3^{2}6^{2}$         & 6  & $24$              & $\phantom{0}2$ & $-2$ & $\phantom{-}0$  & $\phantom{-}2$  & no  & 6 \\
 6B  & $6^{4}$                        & 6  & $24$              & $\phantom{0}0$ & $-4$ & $-2$ & $\phantom{-}0$  & no  & -- \\
 7A  & $1^{3}7^{3}$                   & 7  & $42$              & $\phantom{0}3$ & $-1$ & $-1$ & $\phantom{-}0$  & yes & -- \\
 7B  & $1^{3}7^{3}$                   & 7  & $42$              & $\phantom{0}3$ & $-1$ & $-1$ & $\phantom{-}0$  & yes & -- \\
 8A  & $1^{2}2^{1}4^{1}8^{2}$         & 8  & $16$              & $\phantom{0}2$ & $-2$ & $-2$ & $-2$ & no  & -- \\
 10A & $2^{2}10^{2}$                  & 10 & $20$              & $\phantom{0}0$ & $-4$ & $\phantom{-}0$  & $\phantom{-}2$  & no  & -- \\
 11A & $1^{2}11^{2}$                  & 11 & $11$              & $\phantom{0}2$ & $-2$ & $\phantom{-}2$  & $\phantom{-}0$  & yes & -- \\
 12A & $2^{1}4^{1}6^{1}12^{1}$        & 12 & $12$              & $\phantom{0}0$ & $-4$ & $\phantom{-}0$  & $-2$ & no  & -- \\
 12B & $12^{2}$                       & 12 & $12$              & $\phantom{0}0$ & $-4$ & $\phantom{-}2$  & $\phantom{-}0$  & no  & -- \\
 14A & $1^{1}2^{1}7^{1}14^{1}$        & 14 & $14$              & $\phantom{0}1$ & $-3$ & $\phantom{-}1$  & $\phantom{-}0$  & no  & -- \\
 14B & $1^{1}2^{1}7^{1}14^{1}$        & 14 & $14$              & $\phantom{0}1$ & $-3$ & $\phantom{-}1$  & $\phantom{-}0$  & no  & -- \\
 15A & $1^{1}3^{1}5^{1}15^{1}$        & 15 & $15$              & $\phantom{0}1$ & $-3$ & $\phantom{-}0$  & $-1$ & no  & -- \\
 15B & $1^{1}3^{1}5^{1}15^{1}$        & 15 & $15$              & $\phantom{0}1$ & $-3$ & $\phantom{-}0$  & $-1$ & no  & -- \\
 21A & $3^{1}21^{1}$                  & 21 & $21$              & $\phantom{0}0$ & $-4$ & $-1$ & $\phantom{-}0$  & no  & -- \\
 21B & $3^{1}21^{1}$                  & 21 & $21$              & $\phantom{0}0$ & $-4$ & $-1$ & $\phantom{-}0$  & no  & -- \\
 23A & $1^{1}23^{1}$                  & 23 & $23$              & $\phantom{0}1$ & $-3$ & $-2$ & $\phantom{-}2$  & yes & -- \\
 23B & $1^{1}23^{1}$                  & 23 & $23$              & $\phantom{0}1$ & $-3$ & $-2$ & $\phantom{-}2$  & yes & -- \\
\bottomrule
\end{tabular}
\end{center}
}

\paragraph{Column definitions.}
\begin{itemize}
\item \emph{Frame shape.} Read multiplicatively:
$\prod_i m_i^{e_i}$ with $\sum_i m_i e_i=24$ and $(m_i,e_i)$ the
$g$-orbit lengths on the Golay-code support.
\item \emph{$|\CM(g)|$.} Order of the centraliser in $M_{24}$, ATLAS
column. The anomalous five have centraliser orders
\[
  42,\ 42,\ 11,\ 23,\ 23
\]
for \(7A,7B,11A,23A,23B\).
\item \emph{\(\chi_{24}(g)\).} The number of fixed points in the
permutation representation, equal to the exponent of \(1\) in the
frame shape.
\item \emph{\(c_g(0,0)\).} The \(q^0y^0\)-coefficient of the full K3
twining genus. Since \(c_g(0,1)=c_g(0,-1)=2\), it is
\[
  c_g(0,0)=\chi_{24}(g)-4 .
\]
\item \emph{\(A_g(1),A_g(2)\).} The first massive traces in the
Eguchi--Hikami convention. They are computed from the \(M_{24}\)
character table by
\[
  A_g(1)=\chi_{45}(g)+\chi_{\overline{45}}(g),\qquad
  A_g(2)=\chi_{231}(g)+\chi_{\overline{231}}(g).
\]
\end{itemize}

\paragraph{CHL denominator classification.}
The five primitive CHL denominator rows are
\[
\begin{array}{c|c|c|c}
N & [g] & \text{frame shape} & c_N^{\mathrm{CHL}}(0)\\
\hline
1 & 1A & 1^{24} & 10\\
2 & 2A & 1^{8}2^{8} & 8\\
3 & 3A & 1^{6}3^{6} & 6\\
4 & 4B & 1^{4}2^{2}4^{4} & 4\\
6 & 6A & 1^{2}2^{2}3^{2}6^{2} & 2
\end{array}
\]
Borcherds' weight theorem gives
\[
  \kBKM(\Phi_N)=\frac{c_N^{\mathrm{CHL}}(0)}2,
  \qquad
  \kBKM(\Phi_1,\Phi_2,\Phi_3,\Phi_4,\Phi_6)=(5,4,3,2,1).
\]
Same-order Mathieu classes retain their ordinary twining data in the
26-class table. The CHL column therefore has five entries.

\paragraph{Nikulin boundary constants.}
The order-indexed boundary products at \(N=5,7,8\) carry
\[
\begin{array}{c|c|c|c}
N & [g] & c_N^{\mathrm{Nik}}(0) & c_N^{\mathrm{Nik}}(0)/2\\
\hline
5 & 5A & 4 & 2\\
7 & 7A,7B & 2 & 1\\
8 & 8A & 2 & 1
\end{array}
\]
The elliptic-factor cyclotomic condition \(\varphi(N)\mid 2\) selects
\(N\in\{1,2,3,4,6\}\) for the smooth \(K3\times E\) CHL host.

\paragraph{Cl\'ery--Gritsenko diagonal-divisor constants.}
The eight diagonal-divisor rows are indexed by \((t,N;k)\), with
forms
\[
\begin{array}{c|c|c|c}
\text{form} & (t,N) & \operatorname{wt} & c_{\mathrm{dd}}(0;t,N)\\
\hline
\Delta_5 & (1,1) & 5 & 10\\
\Delta_2 & (2,1) & 2 & 4\\
\nabla_3 & (1,2) & 3 & 6\\
\Delta_1 & (3,1) & 1 & 2\\
\nabla_2 & (1,3) & 2 & 4\\
\Delta_{1/2} & (4,1) & \frac12 & 1\\
\nabla_{3/2} & (1,4) & \frac32 & 3\\
Q_1 & (2,2) & 1 & 2
\end{array}
\]
Thus the row-labelled product constants are
\[
  10,\ 4,\ 6,\ 2,\ 4,\ 1,\ 3,\ 2 .
\]
Only the \(\Delta_5\) row coincides with the primitive CHL \(N=1\) row.

\paragraph{References.}
Eguchi--Hikami Table~1 supplies the frame shapes and permutation
representatives. Gaberdiel--Hohenegger--Volpato Table~2 supplies the
centraliser orders and Table~3 supplies the \(M_{24}\) character values.
Eguchi--Hikami Table~3 supplies the traces \(A_g(n)\). Borcherds'
weight theorem supplies the CHL weights from \(c_N^{\mathrm{CHL}}(0)\).
Cl\'ery--Gritsenko Theorems~1.4 and~3.1 supply the eight
diagonal-divisor rows.
